% ==== Start-up ================================================================
\chapter{Building an AI Venture}
\label{ch:startup}

\section{From an Interesting Model to a Valuable Product}

\begin{tcbitemize}[skin=sectionraster]
\tcbitem[title={Start with a Painful Workflow},raster multicolumn=3]
A start-up begins with a costly, frequent, and observable problem---not with a
model looking for somewhere to live.
Describe who performs the work, what triggers it, which inputs are available,
what outcome matters, and what happens when the decision is wrong.
Customer interviews are evidence about behaviour only when they reconstruct
specific past events; compliments and hypothetical purchase intent are weak
signals\textsuperscript{\cite{constableTalkingHumansSuccess2014}}.

\tcbitem[title={Write the Non-AI Baseline First},raster multicolumn=3]
Before training, compare the proposed system with a rule, search interface,
spreadsheet, process redesign, or human service.
The baseline exposes whether uncertainty-tolerant statistical behaviour is
actually needed.
It also establishes the minimum quality, cost, latency, and reliability that
an AI solution must beat.

\tcbitem[title={Separate Four Hypotheses},raster multicolumn=6]
An AI product simultaneously tests a \keyterm{problem hypothesis} (the pain is
real), a \keyterm{value hypothesis} (improvement changes behaviour), a
\keyterm{feasibility hypothesis} (data and technology meet the target), and a
\keyterm{viability hypothesis} (the economics and obligations are sustainable).
A successful model experiment validates only part of the feasibility
hypothesis.
It does not prove demand, workflow adoption, lawful data rights, or a viable
business.
\end{tcbitemize}

\begin{cautionbox}[The benchmark-to-business fallacy]
An improvement on a public benchmark may have no economic value if the error
being reduced is cheap, the workflow cannot use the output, integration costs
dominate, or the data distribution differs from the customer environment.
Define the operational decision and its error costs before choosing a metric.
\end{cautionbox}

\section{The Evidence Loop}

\begin{figure}[htbp]
\centering
\begin{tikzpicture}[
    node distance=7mm and 9mm,
    venture step/.style={draw,rounded corners=1mm,fill=LimeGreen!8,
        minimum width=2.7cm,minimum height=8mm,align=center,
        font=\sffamily\scriptsize},
    venture arrow/.style={-{Stealth[length=2mm]},thick}]
\node[venture step] (assumption) {Rank the riskiest\\assumption};
\node[venture step,right=of assumption] (test)
    {Design the smallest\\credible test};
\node[venture step,below=of test] (measure)
    {Measure behaviour,\\quality, cost, and risk};
\node[venture step,left=of measure] (decide) {Continue, change,\\or stop};
\draw[venture arrow] (assumption) -- (test);
\draw[venture arrow] (test) -- (measure);
\draw[venture arrow] (measure) -- (decide);
\draw[venture arrow] (decide) -- (assumption);
\end{tikzpicture}
\caption{The start-up evidence loop turns assumptions into explicit tests and
decision gates.}
\label{fig:startup-evidence-loop}
\end{figure}

\begin{tcbitemize}[skin=sectionraster]
\tcbitem[title={Assumption Ledger},raster multicolumn=3]
Maintain one row per important assumption.
Record the consequence of being wrong, present evidence, test method, decision
threshold, owner, and review date.
Rank by \emph{uncertainty multiplied by consequence}, not by which experiment
is most entertaining to build.

\tcbitem[title={Experiment Before Architecture},raster multicolumn=3]
Match the test to the hypothesis.
Use interviews and observation for workflow pain, a concierge service for
value, a frozen representative dataset for technical feasibility, a paid
pilot for willingness to pay, and a threat or compliance review for risks
that cannot ethically be discovered in production.

\tcbitem[title={Evidence Gates},raster multicolumn=6]
Stage gates are useful when they specify evidence rather than ceremony.
Cooper emphasises adaptable, accountable gates with explicit criteria instead
of a rigid linear process\textsuperscript{\cite{cooperPerspectiveStageGateIdeatoLaunch2008}}.
For an AI system, every gate should include product, model, operations, safety,
security, privacy, and economic evidence.
\tcblower
\begin{tblr}{
    width=\linewidth,
    colspec={Q[l,0.16\linewidth]X[l]X[l]X[l]},
    row{1}={font=\bfseries,bg=LimeGreen!15},
    hlines,
    vlines}
Gate & Question & Evidence & Stop or change when \\
Discovery & Is the problem frequent and costly? & Observed workflows,
interviews, current workaround & No repeated behaviour or accountable buyer \\
Feasibility & Can a simple system meet the target? & Baseline, representative
split, error slices, cost profile & Target requires unavailable data or unsafe
behaviour \\
Pilot & Does the workflow improve in context? & Prospective usage, overrides,
time saved, incident log & Users route around the system or errors concentrate \\
Production & Can the service be operated responsibly? & SLOs, monitoring,
security review, legal basis, fallback & Residual risk or unit economics exceed
tolerance \\
Scale & Does value persist across customers? & Cohort retention, support load,
shift tests, margin by segment & Custom work grows faster than repeatable value \\
\end{tblr}
\captionof{table}{Evidence gates for an AI product from discovery to scale.}
\label{tab:startup-evidence-gates}
\end{tcbitemize}

\section{AI Feasibility and Product Architecture}

\begin{tcbitemize}[skin=sectionraster]
\tcbitem[title={Define the Product Contract},raster multicolumn=3]
Specify the input, output, user, decision, latency, availability, acceptable
error, abstention path, privacy boundary, and operating conditions.
Then write the evaluation contract: dataset provenance, split logic, metrics,
slices, uncertainty, baselines, and release threshold.
The product and evaluation contracts must describe the same real-world use.

\tcbitem[title={Build, Buy, or Partner},raster multicolumn=3]
\textbf{Build} when proprietary capability is central and evidence supports
the long-term cost.
\textbf{Buy} when a service is substitutable and contracts meet privacy,
security, availability, and exit needs.
\textbf{Partner} when success requires domain access or distribution that the
team cannot recreate.
Preserve an exit path: exportable data, portable evaluations, versioned
prompts or adapters, and a non-proprietary fallback.

\tcbitem[title={Where Defensibility Comes From},raster multicolumn=6]
Model access alone is rarely a durable moat.
More defensible assets are permissioned data generated by a valuable workflow,
hard-won integration, trust and distribution, a fast learning loop, and
evidence that the system performs safely in a demanding niche.
Feedback data is useful only when its labels are meaningful and its collection
does not create harmful incentives.
\end{tcbitemize}

\begin{engineernote}[Prototype the uncertainty boundary]
A polished interface can hide unreliable behaviour.
During discovery, show confidence limits, abstentions, citations, provenance,
and failure cases explicitly.
The team needs to learn where automation stops being dependable before users
learn that lesson in production.
\end{engineernote}

\section{Unit Economics for AI Systems}

For one completed customer workflow, a useful contribution estimate is
\begin{equation}
  M = P - C_{\mathrm{inference}} - C_{\mathrm{data}}
      - C_{\mathrm{review}} - C_{\mathrm{support}} - C_{\mathrm{risk}},
\end{equation}
where $P$ is realised revenue and the costs include failed requests and human
review, not merely successful model calls.
If a workflow needs $n$ attempts with per-attempt cost $c$ and succeeds with
probability $p$, the expected inference cost per successful workflow is
approximately
\begin{equation}
  C_{\mathrm{inference/success}} = \frac{n c}{p}.
\end{equation}
This approximation must be expanded when retries are correlated, cache hits
matter, long contexts change cost, or expensive failures create support and
liability exposure.

\begin{workedexample}[The cheaper model is not always the cheaper product]
Model A costs \$0.04 per attempt and completes 80\% of workflows without human
review.
Model B costs \$0.10 and completes 95\%.
A failed workflow costs \$1.20 in review time.
Ignoring other costs, expected cost is
$0.04 + 0.20(1.20)=\$0.28$ for A and
$0.10 + 0.05(1.20)=\$0.16$ for B.
Optimise the full workflow, then verify the assumptions with measured queues
and failure categories.
\end{workedexample}

\section{Responsible Commercialisation}

\begin{tcbitemize}[skin=sectionraster]
\tcbitem[title={Data and Rights},raster multicolumn=3]
Record the origin, licence, consent or legal basis, purpose, retention, and
deletion route for each data source.
Separate the right to access data from the right to train a model, generate an
output, or transfer data to a supplier.

\tcbitem[title={Claims and Sales},raster multicolumn=3]
Marketing claims are engineering requirements in public clothing.
Attach scope, date, dataset, uncertainty, and known limits to every material
performance claim.
Do not convert a laboratory result into claims about safety, health, finance,
or legal compliance that the study did not test.

\tcbitem[title={Operating Responsibility},raster multicolumn=3]
Name owners for model changes, security patches, incidents, customer notices,
data-subject requests, and rollback.
Budget for monitoring, red-team work, documentation, insurance, and specialist
review where the domain requires it.

\tcbitem[title={Kill Criteria Are an Asset},raster multicolumn=3]
Stopping preserves time and credibility.
Precommit to kill or pivot criteria for absent demand, unattainable quality,
unsafe residual risk, unlawful data dependence, integration burden, and unit
economics that do not improve with realistic scale.
\end{tcbitemize}

\section{Practical Lab: One-Page AI Product Brief}

\begin{practicallab}[Turn a model idea into an evidence plan]
Choose one workflow and produce a one-page brief containing:
\begin{enumerate}
    \item the user, triggering event, current workaround, and measurable pain;
    \item the non-AI baseline and the smallest AI-assisted intervention;
    \item product and evaluation contracts, including error costs and
          abstention;
    \item the five riskiest assumptions and one cheap, ethical test for each;
    \item a cost model per successful workflow and a sensitivity analysis;
    \item data rights, safety, security, privacy, and legal review questions;
    \item pilot, production, and scale gates with owners and kill criteria.
\end{enumerate}
Interview three relevant users and revise the brief using observed behaviour.
The deliverable is the changed assumption ledger, not a slide deck claiming
that every interview was encouraging.
\end{practicallab}

\begin{checkpoint}
\begin{enumerate}
    \item Which of the four hypotheses does a high offline score support?
    \item What is the strongest non-AI baseline for your proposed workflow?
    \item Which cost grows when model quality falls even if token cost stays
          fixed?
    \item What evidence would make you stop the project next week?
\end{enumerate}
\end{checkpoint}

\section{Further Reading}

Customer interviews and evidence-based discovery are developed in
\cite{constableTalkingHumansSuccess2014}.
Cooper's stage-gate work explains how explicit, adaptable gates support
accountability without forcing innovation into a rigid waterfall
\cite{cooperPerspectiveStageGateIdeatoLaunch2008}.
